% PAGES: 39-45

\section{Diagonal matrices}

\begin{definition}
A square matrix is diagonal if and only if all the entries of the matrix
which are not on the main diagonal are equal to $0$.

In other words, the $n\times n$ matrix $D$ is diagonal if and only if
\begin{equation}
D(j,k) \;=\; 0, \quad \forall\, 1\le j\ne k\le n.
\end{equation}
\end{definition}

Note that, if $d_i=D(i,i)$ for all $i=1:n$, then
\begin{equation}
D \;=\; \begin{pmatrix} d_1 & 0 & \ldots & 0 \\ 0 & d_2 & \ldots & 0 \\
\vdots & \vdots & \ddots & \vdots \\ 0 & 0 & \ldots & d_n \end{pmatrix}.
\end{equation}

From (1.84) and (1.85), it follows that:

(i) any diagonal matrix is symmetric;

(ii) both the sum and the product of two diagonal matrices are diagonal
matrices; see also Lemma 1.11;

(iii) any diagonal matrix is both upper triangular and lower
triangular;\footnote{An equivalent definition for a diagonal matrix is
that a diagonal matrix is a matrix that is both upper triangular and
lower triangular.} see (1.103) and (1.104).

The shorthand notation for the diagonal matrix $D$ given by (1.85) is
\[
D \;=\; \diag(d_k)_{k=1:n}.
\]

The column form of the diagonal matrix $D=\diag(d_k)_{k=1:n}$ is
\begin{equation}
D \;=\; \mathrm{col}\,(d_ke_k)_{k=1:n},
\end{equation}

where $e_k$ is the $k$-th column of the identity matrix of size $n$, for
$k=1:n$.

The row form of the diagonal matrix $D=\diag(d_k)_{k=1:n}$ is
\begin{equation}
D \;=\; \mathrm{row}\,(d_je_j^t)_{j=1:n}.
\end{equation}

\begin{lemma}
Let $A$ be an $m\times n$ matrix with column form
$A=\mathrm{col}\,(a_k)_{k=1:n}$ and row form $A=\mathrm{row}\,(r_j)_{j=1:m}$.

(i) If $D=\diag(d_k)_{k=1:n}$ is a diagonal matrix of size $n$, then
\begin{equation}
AD \;=\; \mathrm{col}\,(d_ka_k)_{k=1:n} \;=\; (d_1a_1\ |\ d_2a_2\ |\ \ldots\ |\ d_na_n).
\end{equation}

(ii) If $D=\diag(d_j)_{j=1:m}$ is a diagonal matrix of size $m$, then
\begin{equation}
DA \;=\; \mathrm{row}\,(d_jr_j)_{j=1:m} \;=\;
\begin{pmatrix} d_1r_1 \\ \text{--} \\ d_2r_2 \\ \text{--} \\ \vdots \\ \text{--} \\ d_nr_n \end{pmatrix}.
\end{equation}

(iii) If $D_1=\diag\!\left(d_j^{(1)}\right)_{j=1:m}$ and
$D_2=\diag\!\left(d_k^{(2)}\right)_{k=1:n}$ are diagonal matrices, then
\begin{equation}
(D_1AD_2)(j,k) \;=\; d_j^{(1)}d_k^{(2)}A(j,k), \quad \forall\, j=1:m,\ k=1:n.
\end{equation}
\end{lemma}

\begin{proof}
(i) Let $D=\mathrm{col}\,(d_ke_k)_{k=1:n}$ be the column form of the
diagonal matrix $D$, where $e_k$ is the $k$-th column of the identity
matrix of size $n$. Recall from (1.26) that $Ae_k=a_k$, for all $k=1:n$.
Then, using (1.11), we find that
\[
AD \;=\; \mathrm{col}\,(A\cdot d_ke_k)_{k=1:n} \;=\; \mathrm{col}\,(d_k\cdot Ae_k)_{k=1:n}
\;=\; \mathrm{col}\,(d_ka_k)_{k=1:n}.
\]

(ii) Let $D=\mathrm{row}\,(d_je_j^t)_{j=1:m}$ be the row form of the
diagonal matrix $D$, where $e_j$ is the $j$-th column of the identity
matrix of size $m$. Recall from (1.27) that $e_j^tA=r_j$, for all $j=1:m$.
Then, using (1.12), we find that
\[
DA \;=\; \mathrm{row}\,(d_je_j^t\cdot A)_{j=1:m} \;=\; \mathrm{row}\,(d_j\cdot e_j^tA)_{j=1:m}
\;=\; \mathrm{row}\,(d_jr_j)_{j=1:m}.
\]

(iii) Let
\[
D_1 \;=\; \mathrm{row}\,\!\left(d_j^{(1)}(e_j^{(1)})^t\right)_{j=1:m}
\quad\text{and}\quad
D_2 \;=\; \mathrm{col}\,\!\left(d_k^{(2)}e_k^{(2)}\right)_{k=1:n}
\]
be the row form of the diagonal matrix $D_1$ and the column form of the
diagonal matrix $D_2$, respectively, where $e_j^{(1)}$ is the $j$-th
column of the identity matrix of size $m$ and $e_k^{(2)}$ is the $k$-th
column of the identity matrix of size $n$.

Let $1\le j\le m$ and $1\le k\le n$ arbitrary. Recall from (1.14) that the
entry $(j,k)$ of the matrix $D_1AD_2$ is
\begin{equation}
(D_1AD_2)(j,k) \;=\; \left(d_j^{(1)}(e_j^{(1)})^t\right)A\left(d_k^{(2)}e_k^{(2)}\right)
\;=\; d_j^{(1)}d_k^{(2)}\left((e_j^{(1)})^tAe_k^{(2)}\right).
\end{equation}

Recall from (1.26) that $Ae_k^{(2)}=a_k$. Since $a_k=(A(i,k))_{i=1:m}$, we
find using (1.5) that
\begin{equation}
(e_j^{(1)})^tAe_k^{(2)} \;=\; (e_j^{(1)})^ta_k \;=\; \sum_{i=1}^m e_j^{(1)}(i)A(i,k) \;=\; A(j,k),
\end{equation}
since $e_j^{(1)}(i)=0$ for any $i\ne j$ and $e_j^{(1)}(j)=1$.

From (1.91) and (1.92), we conclude that
\[
(D_1AD_2)(j,k) \;=\; d_j^{(1)}d_k^{(2)}A(j,k), \quad \forall\, j=1:m,\ k=1:n. \qedhere
\]
\end{proof}

\begin{lemma}
The product of two diagonal matrices of the same size is a diagonal
matrix. Also, products of diagonal matrices are commutative.

In other words, if $D_1=\diag\!\left(d_k^{(1)}\right)_{k=1:n}$ and
$D_2=\diag\!\left(d_k^{(2)}\right)_{k=1:n}$ are diagonal matrices, then
\begin{equation}
D_1D_2 \;=\; \diag\!\left(d_k^{(1)}d_k^{(2)}\right)_{k=1:n}
\end{equation}
and $D_1D_2=D_2D_1$.
\end{lemma}

\begin{proof}
Let $D_1=\diag\!\left(d_k^{(1)}\right)_{k=1:n}$ and
$D_2=\diag\!\left(d_k^{(2)}\right)_{k=1:n}$ be diagonal matrices.

The column form of $D_1$ is $D_1=\mathrm{col}\,\!\left(d_k^{(1)}e_k\right)_{k=1:n}$,
and, using (1.88), we find that
\[
D_1D_2 \;=\; \mathrm{col}\,\!\left(d_k^{(2)}\cdot d_k^{(1)}e_k\right)_{k=1:n}
\;=\; \mathrm{col}\,\!\left(d_k^{(1)}d_k^{(2)}e_k\right)_{k=1:n}
\;=\; \diag\!\left(d_k^{(1)}d_k^{(2)}\right)_{k=1:n}.
\]

By using (1.93) for $D_1=D_2$ and $D_2=D_1$, we find that
\[
D_2D_1 \;=\; \diag\!\left(d_k^{(2)}d_k^{(1)}\right)_{k=1:n} \;=\; \diag\!\left(d_k^{(1)}d_k^{(2)}\right)_{k=1:n} \;=\; D_1D_2,
\]
and we conclude that products of diagonal matrices are commutative.
\end{proof}

\begin{lemma}
(i) A diagonal matrix is nonsingular if and only if all its entries on the
main diagonal are nonzero.

In other words, the matrix $D=\diag(d_k)_{k=1:n}$ is nonsingular if and
only if $d_k\ne 0$ for all $k=1:n$.

(ii) If $D=\diag(d_k)_{k=1:n}$ is a nonsingular diagonal matrix, then
\begin{equation}
D^{-1} \;=\; \diag\!\left(\frac{1}{d_k}\right)_{k=1:n}.
\end{equation}
\end{lemma}

\begin{proof}
(i) Let $D=\diag(d_k)_{k=1:n}$ be a diagonal matrix. Then,
$\det(D)=\prod_{k=1}^n d_k$; cf.\ (10.1). From Theorem 1.2, we conclude
that the matrix $D$ is nonsingular if and only if $d_k\ne 0$ for all
$k=1:n$.

(ii) Let $D^{-1}$ be the matrix given by (1.94); recall that $d_k\ne 0$
for all $k=1:n$, since $D$ is nonsingular. From (1.93), we find that
\[
DD^{-1} \;=\; \diag\!\left(d_k\cdot\frac{1}{d_k}\right)_{k=1:n} \;=\; \diag(1)_{k=1:n} \;=\; I,
\]
and therefore $D^{-1}$ is the inverse matrix of $D$; cf.\ Lemma 1.6.
\end{proof}

\subsection{Converting between covariance and correlation matrices}

The covariance matrix $\Sigma_{\bfX}$ and the correlation matrix
$\Omega_{\bfX}$ of $n$ nonconstant random variables $X_1,X_2,\ldots,X_n$
are the $n\times n$ matrices given by
\begin{align*}
\Sigma_{\bfX}(j,k) &= \cov(X_j,X_k), \quad \forall\, 1\le j,k\le n; \\
\Omega_{\bfX}(j,k) &= \corr(X_j,X_k), \quad \forall\, 1\le j,k\le n,
\end{align*}
where $\cov(X_j,X_k)$ and $\corr(X_j,X_k)$ denote the covariance and the
correlation of the random variables $X_j$ and $X_k$, respectively. Recall
that
\[
\cov(X_j,X_k) \;=\; \sigma_j\sigma_k\corr(X_j,X_k), \quad \forall\, 1\le j,k\le n,
\]
where $\sigma_j$ and $\sigma_k$ are the standard deviations of $X_j$ and
$X_k$, respectively. Thus,
\[
\Sigma_{\bfX}(j,k) \;=\; \sigma_j\sigma_k\corr(X_j,X_k), \quad \forall\, 1\le j,k\le n.
\]

Then, from (1.90), it follows that
\begin{equation}
\Sigma_{\bfX} \;=\; D_{\sigma_{\bfX}}\Omega_{\bfX}D_{\sigma_{\bfX}},
\end{equation}
where $D_{\sigma_{\bfX}}$ is the diagonal matrix given by
$D_{\sigma_{\bfX}}=\diag(\sigma_i)_{i=1:n}$.

We conclude that, if the standard deviations of the $n$ random variables
are known, then the covariance matrix $\Sigma_{\bfX}$ can be uniquely
determined from the correlation matrix $\Omega_{\bfX}$ by using formula
(1.95).

Since $\sigma_i\ne 0$ for all $i=1:n$, the matrix $D_{\sigma_{\bfX}}$ is
nonsingular and
$(D_{\sigma_{\bfX}})^{-1}=\diag\!\left(\frac{1}{\sigma_i}\right)_{i=1:n}$;
see (1.94). Then, from (1.95), we obtain that
\begin{equation}
\Omega_{\bfX} \;=\; (D_{\sigma_{\bfX}})^{-1}\Sigma_{\bfX}(D_{\sigma_{\bfX}})^{-1}.
\end{equation}

Note that
\begin{equation}
D_{\sigma_{\bfX}} \;=\; \diag(\sigma_i)_{i=1:n} \;=\; \diag\!\left(\sqrt{\Sigma_{\bfX}(i,i)}\right)_{i=1:n},
\end{equation}
since, by definition, $\sigma_i^2=\cov(X_i,X_i)=\Sigma_{\bfX}(i,i)$. In
other words, the entries of the matrix $D_{\sigma_{\bfX}}$ can be
obtained from the entries of $\Sigma_{\bfX}$.

We conclude that the correlation matrix $\Omega_{\bfX}$ is uniquely
determined by the covariance matrix $\Sigma_{\bfX}$, and can be computed
by using (1.96) and (1.97).

For more properties of covariance and correlation matrices, see section
7.1.

\begin{examplex}
Assume that the covariance matrix of three random variables is
\[
\Sigma_{\bfX} \;=\; \begin{pmatrix} 1 & 0.1 & -0.6 \\ 0.1 & 0.25 & 0.5 \\ -0.6 & 0.5 & 4 \end{pmatrix}.
\]

Note that
\[
\sigma_1=\sqrt{\Sigma_{\bfX}(1,1)}=1; \quad \sigma_2=\sqrt{\Sigma_{\bfX}(2,2)}=0.5; \quad \sigma_3=\sqrt{\Sigma_{\bfX}(3,3)}=2.
\]

Thus, $D_{\sigma_{\bfX}}=\begin{pmatrix} 1 & 0 & 0 \\ 0 & \frac12 & 0 \\ 0 & 0 & 2 \end{pmatrix}$
and $(D_{\sigma_{\bfX}})^{-1}=\begin{pmatrix} 1 & 0 & 0 \\ 0 & 2 & 0 \\ 0 & 0 & \frac12 \end{pmatrix}$.
From (1.96), we find that the correlation matrix of the three random
variables is
\[
\Omega_{\bfX} \;=\; (D_{\sigma_{\bfX}})^{-1}\Sigma_{\bfX}(D_{\sigma_{\bfX}})^{-1}
\;=\; \begin{pmatrix} 1 & 0.2 & -0.3 \\ 0.2 & 1 & 0.5 \\ -0.3 & 0.5 & 1 \end{pmatrix}.
\]
\end{examplex}

\begin{examplex}
Assume that the correlation matrix of three random variables with
standard deviations $\sigma_1=2$, $\sigma_2=1$, and $\sigma_3=0.9$, is
\[
\Omega_{\bfX} \;=\; \begin{pmatrix} 1 & 0.15 & -0.2 \\ 0.15 & 1 & 0.1 \\ -0.2 & 0.1 & 1 \end{pmatrix}.
\]

Then, $D_{\sigma_{\bfX}}=\begin{pmatrix} 2 & 0 & 0 \\ 0 & 1 & 0 \\ 0 & 0 & 0.9 \end{pmatrix}$,
and, from (1.95), we conclude that the covariance matrix of the three
random variables is
\[
\Sigma_{\bfX} \;=\; D_{\sigma_{\bfX}}\Omega_{\bfX}D_{\sigma_{\bfX}}
\;=\; \begin{pmatrix} 4 & 0.3 & -0.36 \\ 0.3 & 1 & 0.09 \\ -0.36 & 0.09 & 0.81 \end{pmatrix}.
\]
\end{examplex}

Formulas similar to (1.95) and (1.96) can also be derived for sample
covariance matrices and sample correlation matrices; see (1.99) and
(1.100), respectively.

Denote by $\widehat{\sigma}_i=\sqrt{\widehat{\var}(X_i)}$ the sample
standard deviation of $X_i$, for $i=1:n$. Then, the sample correlation of
$X_j$ and $X_k$ is
\[
\widehat{\corr}(X_j,X_k) \;=\; \frac{\widehat{\cov}(X_j,X_k)}{\widehat{\sigma}_j\widehat{\sigma}_k}, \quad \forall\, 1\le j,k\le n,
\]
where $\widehat{\cov}(X_j,X_k)$ is the sample covariance of $X_j$ and
$X_k$ given by (1.31).

The sample covariance matrix $\widehat{\Sigma}_{\bfX}$ and the sample
correlation matrix $\widehat{\Omega}_{\bfX}$ of the random variables
$X_1,X_2,\ldots X_n$ are the $n\times n$ matrices given by
\begin{align*}
\widehat{\Sigma}_{\bfX}(j,k) &= \widehat{\cov}(X_j,X_k), \quad \forall\, 1\le j,k\le n; \\
\widehat{\Omega}_{\bfX}(j,k) &= \widehat{\corr}(X_j,X_k), \quad \forall\, 1\le j,k\le n.
\end{align*}

If
\begin{equation}
\widehat{D}_{\sigma_{\bfX}} \;=\; \diag(\widehat{\sigma}_i)_{i=1:n} \;=\; \diag\!\left(\sqrt{\widehat{\Sigma}_{\bfX}(i,i)}\right)_{i=1:n},
\end{equation}
then,
\begin{align}
\widehat{\Sigma}_{\bfX} &= \widehat{D}_{\sigma_{\bfX}}\widehat{\Omega}_{\bfX}\widehat{D}_{\sigma_{\bfX}}; \\
\widehat{\Omega}_{\bfX} &= (\widehat{D}_{\sigma_{\bfX}})^{-1}\widehat{\Sigma}_{\bfX}(\widehat{D}_{\sigma_{\bfX}})^{-1}.
\end{align}

\par\medskip\noindent\textit{Example (continued):}\
The sample covariance matrix $\widehat{\Sigma}_{\bfX}$ of the daily
returns of AAPL, FB, GOOG, MSFT, YHOO between 1/11/2013 and 1/29/2013 is
given by (1.44). The diagonal matrix of the standard deviation of the
daily returns
$\widehat{D}_{\sigma_{\bfX}}=\diag\!\left(\sqrt{\widehat{\Sigma}_{\bfX}(i,i)}\right)_{i=1:5}$
is
\begin{equation}
\widehat{D}_{\sigma_{\bfX}} \;=\;
\begin{pmatrix}
0.0426 & 0 & 0 & 0 & 0 \\
0 & 0.0254 & 0 & 0 & 0 \\
0 & 0 & 0.0194 & 0 & 0 \\
0 & 0 & 0 & 0.0072 & 0 \\
0 & 0 & 0 & 0 & 0.0147
\end{pmatrix}.
\end{equation}

From (1.100), and using (1.44) and (1.101), we find that the sample
correlation matrix $\widehat{\Omega}_{\bfX}$ of the daily returns of
AAPL, FB, GOOG, MSFT, YHOO between 1/11/2013 and 1/29/2013
is\footnote{Note that, when computing $\widehat{D}_{\sigma_{\bfX}}$ and
$\widehat{\Omega}_{\bfX}$ from (1.101) and (1.102), respectively, we
used the non-truncated values for the entries of the sample covariance
matrix $\widehat{\Sigma}_{\bfX}$, and not the values from (1.44) which
are rounded at four decimal digits.}
\begin{equation}
\widehat{\Omega}_{\bfX} \;=\;
\begin{pmatrix}
1 & -0.0087 & -0.0736 & -0.0615 & -0.2002 \\
-0.0087 & 1 & 0.1075 & -0.0506 & 0.2929 \\
-0.0736 & 0.1075 & 1 & 0.5614 & 0.0591 \\
-0.0615 & -0.0506 & 0.5614 & 1 & 0.0250 \\
-0.2002 & 0.2929 & 0.0591 & 0.0250 & 1
\end{pmatrix}.
\end{equation}
\hfill$\square$\par\medskip
